\section{Discussion}
\label{sec:discussion}

\subsection{The Boundary of Compositionality Is the Boundary of Linearity}

Our results point to a surprisingly clean answer to the question ``which linear directions are compositional?''
The directions that form coherent linear representations compose reliably; those that do not form coherent directions cannot be composed at all.
The real divide is not between ``compositional'' and ``non-compositional'' linear directions --- it is between concepts that are linearly represented and concepts that are not.

Morphological transformations (verb tenses, pluralization, derivational morphology) are the prototypical compositional directions.
They encode regular, functional transformations with consistent behavior across word pairs, producing \loo scores above~0.3 and near-perfect steering composition.
Encyclopedic associations (country$\to$capital, thing$\to$color) are moderately linear and compose adequately.
Lexicographic semantic relations (synonymy, antonymy, meronymy, hyponymy) fail because ``synonym'' is not a single transformation --- the direction from ``big'' to ``large'' is fundamentally different from the direction from ``fast'' to ``quick.''
These relations are better understood as \emph{family-resemblance categories} that lack a unifying direction.

\subsection{Why Lexicographic Relations Fail}

Three factors likely explain the near-zero \loo consistency of lexicographic relations.
First, these are not \emph{single transformations}: unlike ``make plural,'' which applies a consistent morphological operation, ``find a synonym'' requires different operations for each word.
Second, \emph{superposition}: as \citet{elhage2022superposition} show, infrequent or complex features may be encoded in superposed, distributed representations rather than clean directions.
Lexicographic relations span diverse semantic fields and are unlikely to be compressed into a single direction.
Third, \emph{context-dependence}: \citet{opielka2026causality} demonstrate that the extracted direction for a concept can depend on the format of extraction.
Lexicographic relations may be especially sensitive to this effect because they are defined relationally rather than functionally.

\subsection{Two Senses of Compositionality}

Our metrics reveal a tension between two senses of ``compositionality'' that the literature has not clearly distinguished:

\para{Geometric composition.}
The sum $\vd_A + \vd_B$ preserves both component directions as measured by \cfs.
This is \emph{easier} for aligned directions (within-category): when $\vd_A \cdot \vd_B > 0$, the sum is closer to each component than when the directions are orthogonal, producing higher \cfs.
However, alignment also means that each direction partially activates the other, increasing interference.

\para{Functional composition.}
Steering with $\vd_A + \vd_B$ achieves both intended effects without degrading either.
This works well for both aligned and orthogonal directions, as long as both individual directions are coherent.
The practical implication is clear: {\bf steering vector composition works in the regime where individual directions are well-formed (\loo $> 0.2$), regardless of whether the directions are aligned or orthogonal.}

\subsection{Connection to Prior Work}

Our causal inner product analysis confirms the block-diagonal structure reported by \citet{park2024linear}, extending their finding to show that this structure holds for morphological and encyclopedic concepts but degrades for lexicographic relations.
The finding that ``some function vectors compose and some don't''~\citep{todd2024function} aligns with our taxonomy: the compositions that work involve well-defined functional transformations (our morphological categories), while those that fail likely involve poorly defined directions.
\citet{elhage2022superposition}'s superposition framework explains why lexicographic relations do not form clean directions --- they are encoded in distributed, superposed representations that resist extraction into single vectors.

\subsection{Limitations}
\label{sec:limitations}

Our study has several limitations that constrain the scope of our conclusions.

\para{Single model.}
We test only \pythia.
Larger models may resolve superposition for some concepts~\citep{elhage2022superposition}, potentially making lexicographic relations more linear.
Architecturally different models (mixture-of-experts, state-space models) may organize representations differently.

\para{Unembedding space only.}
We extract directions from the unembedding matrix rather than intermediate layers.
Composition structure may differ across layers, and intermediate representations may capture concepts that the unembedding space does not.

\para{Limited concept coverage.}
\bats covers morphological and basic semantic relations but not the abstract behavioral concepts (honesty, sycophancy) tested in the steering literature~\citep{zou2023representation,panickssery2023steering}.
Whether our findings extend to behavioral concepts remains an open question.

\para{Ceiling effects in steering.}
Functional steering accuracy exceeds 96\% everywhere, providing little discriminative power between pair types.
Harder steering tasks or finer-grained metrics (logit shift magnitude, rank changes) would enable more nuanced comparisons.

\para{Single-token constraint.}
Requiring both words to be single tokens limits the available pairs and may introduce selection bias toward shorter, more common words.

\para{No ground truth.}
We define compositionality operationally through \cfs and steering accuracy.
There is no established ground-truth benchmark for direction compositionality, making it difficult to validate our metrics independently.
